\subsection{Sensitivity analysis} \label{sec-results-sensitivity}

In this section we assess the sensitivity of the model outcomes to its parameters $C_{L\max}$, $C_{D0}$, and $C_{D\max}$ based on the initial speed required to achieve a touchdown speed of 0.1 and 0.4 (Fig.~\ref{fig-unsteady}). We increase $C_{L\max}$ and $C_{D\max}$ well beyond their steady-state values as a preliminary exploration of how unsteady aerodynamic effects (which are not explicitly captured in our aerodynamics model) may affect the landing. The main unsteady aerodynamic effect of interest is dynamic or delayed stall due to the rate of change of angle of attack, which can result in large, transient increases in the lift and drag coefficients \cite{stricklandForceCoefficientsNACA00151987,granlundUnsteadyPitchingFlat2013,leishmanPrinciplesHelicopterAerodynamics2017,kieferDynamicStallHigh2022}. Previous work has found that $C_{L\max}$ and $C_{D\max}$ scale linearly with pitch rate, whilst approximately retaining the trigonometric form of Eq.~\ref{eq-clcd} \cite{stricklandForceCoefficientsNACA00151987,granlundUnsteadyPitchingFlat2013}, supporting the approach taken here. Transient increase in lift and drag due to added mass effects \cite{andersenUnsteadyAerodynamicsFluttering2005,wangUnsteadyForcesFlows2004} can also be accounted for in this way, but phase lag effects are not.

For a small touchdown speed of 0.1, doubling $C_{L\max}$ from its nominal value resulted in a 44\% reduction in the initial speed required and a 65\% and 55\% reduction in the horizontal and vertical displacements, respectively (grey and blue trajectories in Fig.~\ref{fig-unsteady}A). However, doubling $C_{D\max}$ from its nominal value led to a 31\% increase in initial speed and a 68\% and 15\% increase in the horizontal and vertical displacements (grey and yellow trajectories in Fig.~\ref{fig-unsteady}A). Increasing $C_{D\max}$ led to increases in the initial speed because low speeds can only be achieved with steep climb angles. Increasing $C_{D\max}$ reduces the ability to achieve these steep climbs as it increases the deceleration relative to the increase in flight path angle (see Section~\ref{sec-results-optcontrol}) and as such requires a greater initial speed. For larger touchdown speeds of 0.4 this effect was less pronounced and doubling $C_{D\max}$ decreased the required initial speed (5\%), although this was still much less than when doubling $C_{L\max}$ (49\%).

The largest initial speeds are capped at the maximum terminal speed,
\[
    \nu_T^* = \sqrt{\frac{1}{C_{D0}}},
\]
in Fig.~\ref{fig-unsteady}. This is the fastest speed a gliding bird could travel at without additional energy input (i.e. flapping the wings). As such, regions in the parameter space that require an initial speed greater than $\nu_T^*$ would be unable to attain the given touchdown speed---even in unconstrained flight. Such limits could be important in understanding the morphological requirements early in the evolution of flight.

Birds \cite{carruthersAutomaticAeroelasticDevices2007,poletRapidAreaChange2015,kleinheerenbrinkOptimizationAvianPerching2022} and other gliding vertebrates \cite{bahlmanGlidePerformanceAerodynamics2013,khandelwalHowBiomechanicsPath2020} have been observed to rapidly pitch up prior to landing which can result in large increases in their lift and drag coefficients \cite{granlundUnsteadyPitchingFlat2013,poletUnsteadyDynamicsRapid2015}. This analysis shows how such kinematics (and the resulting aerodynamic effects) could be of benefit when landing. Large increases in $C_{L\max}$ reduce the initial speed required to land at a given touchdown speed (Fig.~\ref{fig-unsteady}) and thus would require less energy input, either through flapping the wings to accelerate or diving from higher heights. It also reduces the horizontal and vertical displacements covered during the landing manoeuvre. Together these effects could make unpowered landings more feasible within the environments birds fly.

While unsteady aerodynamic effects will not affect the minimum drag coefficient, $C_{D0}$, as it relates mainly to skin-friction drag, we assess the sensitivity of the model to this parameter as well. For a touchdown speed of 0.1, doubling and halving $C_{D0}$ from its nominal value (0.03) resulted in a 3\% larger and 2\% smaller required initial speed, respectively.

\begin{figure}[!htbp]
    \centering
    \includegraphics[width=\textwidth]{sensitivity}
    \caption{Sensitivity of the model to its parameters $C_{L\max}$ and $C_{D\max}$. Sensitivity is assessed by the required initial speed needed to reach a touchdown speed of 0.1 (A) and 0.4 (B) for the given set of parameters. Physical-space trajectories are shown below for example parameter sets indicated in the parameter space. The grey trajectory corresponds to the values for $C_{L\max}$ and $C_{D\max}$ that have been used throughout.}
    \label{fig-unsteady}
\end{figure}
